\section{Related Work}
\label{sec:related}

\paragraph{Positive-definite matrix geometry.}
The affine-invariant geometry of positive-definite matrices provides the
distance, exponential map, congruence invariance, and nonpositive curvature
used throughout
\cite{bhatia2007positive,pennec2006riemannian,higham2008functions}.
Geodesically convex optimization and best approximation on SPD manifolds are
developed in
\cite{sra2015conic,lim2004best,bacak2014convex,tumpach2024totally}.
Spectral subdifferential calculus follows the Hermitian convex-analysis
framework of \cite{lewis1996convex}.  Our contribution is the joint
specialization to structured condition targets, explicit penalty and
certificate interfaces, and the exact comparison between metric projection
and condition-number optimization.

\paragraph{Scaling and structured preconditioning.}
Diagonal and block condition-number optimization is connected to optimal
scaling, equilibration, and semidefinite formulations
\cite{bauer1963optimally,vandersluis1969equilibration,
shapiro1982optimally,shapiro1985block,marechal2009optimizing,
lu2011minimizing,jambulapati2020fast,qu2024optimal,gao2023scalable,
ghadimi2025omega}.
In particular, mixed packing--covering SDP solvers and customized large-scale
methods address the computational diagonal problem.  Geodesically convex
preconditioning over symmetric Lie-group quotients provides general
condition-number convexity, projected gradients, and first-order complexity
bounds \cite{dogan2025geodesicpreconditioning}.  We use that established
convexity viewpoint as background.  Applied to a Kronecker transformation
group, that framework supplies condition-objective convexity and first-order
optimization, but does not compare the optimizer with the AIRM projection or
derive the extreme-state marginal criterion.  Our diagonal and block results retain the
exact Loewner reachability question in the notation used later for Kronecker
target distance; the new comparison is between AIRM projection and spectral
condition optimality on the Kronecker SPD submanifold.

Diagonal adaptive methods include AdaGrad and Adam
\cite{duchi2011adagrad,kingma2014adam}; Adafactor and SM3 use factored or
shared statistics \cite{shazeer2018adafactor,anil2019sm3}; K-FAC and KFC use
Kronecker curvature approximations \cite{martens2015kfac,grosse2016kfc};
and Shampoo uses tensor-factor preconditioners
\cite{gupta2018shampoo,anil2020scalable}.  These algorithms motivate the
matrix families.  Identifying a concrete optimizer state with a point in the
family additionally requires a realization theorem specifying damping, bias
correction, update scale, and interpolation.

\paragraph{Kronecker approximation and covariance geometry.}
Kronecker approximation is classical under Frobenius loss
\cite{vanloan1993kronecker}.  Application-driven nearest-Kronecker
preconditioners use Frobenius or operator-specific product approximations for
Toeplitz, imaging, Markov/SAN, and stochastic Galerkin systems
\cite{kamm2000optimal,nagy2006imaging,langville2004san,
langville2004testing,ullmann2010stochastic}.  Those works ask whether a
computationally convenient product approximates a structured operator well or
accelerates a particular solver.  Our comparison instead fixes one Kronecker
family and separates its intrinsic AIRM projection from its exact
Hessian-relative condition-number minimizer.

Separable covariance and matrix-normal models
lead to likelihood and flip-flop procedures
\cite{dutilleul1999mle,werner2008kronecker}.  Geodesic convexity for
covariance estimation and Kronecker structure is established in
\cite{wiesel2012geodesic,wiesel2012kroneckerconvexity}; a broader survey is
\cite{wang2022kroneckersurvey}.  The Fisher/AIRM
quotient geometry, determinant-one factorization, scale gauge, and
dimension-weighted product metric are developed in
\cite{mccormack2025kroneckergeometry}, with related Riemannian algorithms in
\cite{bouchard2021onlinekronecker,simonis2025geodesicvb}.  Work connecting
Shampoo factors to Kronecker approximation studies a complementary algorithmic
question \cite{morwani2024shampoo}.

\begin{center}
\small
\begin{tabular}{p{0.24\linewidth}p{0.31\linewidth}p{0.35\linewidth}}
\toprule
Problem & Established object & Role in this paper \\
\midrule
Kronecker approximation & Frobenius nearest product
\cite{vanloan1993kronecker} & Contrast with intrinsic AIRM projection \\
Separable covariance & Likelihood equations and geodesic convexity
\cite{dutilleul1999mle,wiesel2012kroneckerconvexity} & Distinguish statistical
likelihood from metric projection \\
Kronecker information geometry & Scale quotient and product Fisher/AIRM
metric \cite{mccormack2025kroneckergeometry} & Infrastructure for logarithmic
partial-trace projection residuals \\
Optimal preconditioning & Geodesically convex condition optimization and
first-order methods \cite{dogan2025geodesicpreconditioning} & Exact criterion
for when the AIRM projection belongs to the condition-number argmin \\
\bottomrule
\end{tabular}
\end{center}

Thus the quotient metric itself is supporting infrastructure.  The main
increment here is the global AIRM projection characterization, its residual
error and reachability brackets, the extreme-state marginal criterion for
projection optimality, and strict separation from best Hessian-relative
preconditioning.
